% ====== Tipo de archivo ======
%\documentclass[journal]{IEEEtran}
\documentclass{article}

% ====== Paquetes básicos ======
\usepackage[utf8]{inputenc}    % Codificación UTF-8
\usepackage[T1]{fontenc}       % Acentos correctos
\usepackage[spanish]{babel}    % Traducción al español
\usepackage{csquotes}          % Recomendado para biblatex con babel
\usepackage{hyperref}          % Hipervínculos
\usepackage{amsmath,amssymb}   % Matemáticas
\usepackage{graphicx}          % Imágenes
\usepackage{tikz}              % Diagramas vectoriales
\usetikzlibrary{arrows.meta,positioning}
\usepackage{caption}           % Mejor control de captions
\usepackage{geometry}          % Márgenes
\usepackage{fancyhdr}          % Encabezados y pies
\usepackage{parskip}           % Mejor separación entre párrafos

% ====== Referencias ======
\usepackage[
    backend=biber,
    style=numeric-comp,     % cambie a APA de ser necesario
    sorting=none
]{biblatex} 
\bibliography{references}

% agrege `\printbibliography{}` para las referencias
% cite con `\cite{}`
% agregue las referencias a un `references.bib`

% ====== Configuración ======
\geometry{margin=2.5cm}
\pagestyle{fancy}
\fancyhf{}
\rhead{\thepage}
\setlength{\headheight}{13pt}
\lhead{\leftmark}

% ====== Comando personalizado para imágenes ======
% Uso: \imagen{ruta}{Texto del pie de figura}
\newcommand{\imagen}[2]{
    \begin{figure}[h!]
        \centering
        \includegraphics[width=200px]{#1}
        \caption{#2}
    \end{figure}
}

% ====== Datos del artículo ======
\title{Portada}
\author{Lester Iván Hernández Hernandez}


\usepackage[spanish]{babel}
\usepackage{datetime}
\newdateformat{mydateformat}{\THEDAY/\THEMONTH/\THEYEAR}

\newcommand{\titlename}{Asignación 2do Parcial}
\newcommand{\signature}{Clase (QQ-100)}
\newcommand{\seccion}{1100}
\newcommand{\docente}{Alex Conrado Mejía Valladares}
%Individual
\newcommand{\alumno}{Lester Iván Hernández Hernandez}
\newcommand{\cuenta}{20242000637}

%Grupo
%\newcommand{\grupo}{2}
%\newcommand{\integrantes}{}

\newcommand{\portada}{
    \begin{titlepage}
        \centering
        \vspace*{-1cm}
        \includegraphics[width=4.5cm]{Logo-unah.jpg}\\
        \vspace*{0.5cm}
        {\large\bfseries UNIVERSIDAD NACIONAL AUTÓNOMA DE HONDURAS
        \\UNAH-VS\par}
        \vspace{2cm}
        {\large \titlename\par}
        \vspace{1cm}
        {\large Asignatura: \signature\par}
        {\large Sección: \seccion\par}
        {\large Docente: \docente\par}
        \vspace{1cm}
        %Individual
        {\large Alumno: \alumno\par}
        {\large N. de cuenta: \cuenta\par}
        %Grupo
        %{\large Grupo: \grupo\par}
        %{\large Integrantes: \integrantes\par}
        \vfill
        {\large \today\par}
        \vspace*{1in}
    \end{titlepage}
}


\begin{document}
    \portada{}
    \newpage
    
    \section{Proceso de Electrodeposición de Metales}

Un proceso industrial que integra de manera natural los tres tipos de compuestos solicitados es el \textbf{recubrimiento electrolítico de metales} (por ejemplo, niquelado o cromado de una pieza de acero). Se emplea para proteger piezas metálicas de la corrosión o para darles un acabado decorativo. A continuación se describen sus etapas.

\subsection{Desengrase alcalino (uso de una base)}

Antes de recubrir la pieza, es necesario eliminar grasas, aceites y residuos de manufactura de su superficie.

\begin{itemize}
    \item Se sumerge la pieza en una solución alcalina caliente, típicamente hidróxido de sodio (NaOH) o carbonato de sodio (Na$_2$CO$_3$).
    \item La base saponifica las grasas (las convierte en jabones solubles en agua) mediante hidrólisis alcalina, dejando la superficie metálica limpia y libre de compuestos orgánicos.
\end{itemize}

Reacción general (con una grasa tipo éster):
\begin{equation}
    \text{Grasa (éster)} + \text{NaOH} \rightarrow \text{Jabón (sal soluble)} + \text{Glicerina}
\end{equation}

\subsection{Decapado o activado (uso de un ácido)}

Después del enjuague, la pieza pasa por un baño ácido.

\begin{itemize}
    \item Se emplea ácido clorhídrico (HCl) o ácido sulfúrico (H$_2$SO$_4$) diluido.
    \item Su función es disolver la capa de óxido metálico (herrumbre u óxidos naturales) formada sobre la superficie, dejando el metal ``activado'' y listo para que el recubrimiento se adhiera correctamente.
\end{itemize}

Reacción representativa (óxido de hierro):
\begin{equation}
    \text{Fe}_2\text{O}_3 + 6\text{HCl} \rightarrow 2\text{FeCl}_3 + 3\text{H}_2\text{O}
\end{equation}

\subsection{Enjuague intermedio}

Entre cada etapa se enjuaga la pieza con agua para evitar que los reactivos de un baño contaminen el siguiente.

\subsection{Electrodeposición (uso de una sal metálica)}

Esta es la etapa central del proceso.

\begin{itemize}
    \item La pieza (cátodo) se sumerge en una celda electrolítica que contiene una sal metálica disuelta, como sulfato de níquel (NiSO$_4\cdot$6H$_2$O) para niquelado, o sulfato de cobre (CuSO$_4$) para cobreado.
    \item Un ánodo del mismo metal (por ejemplo, níquel puro) se conecta al polo positivo de una fuente de corriente directa.
\end{itemize}

Al pasar la corriente:

En el cátodo (la pieza), los iones metálicos se reducen y se depositan como metal sólido:
\begin{equation}
    \text{Ni}^{2+} + 2e^- \rightarrow \text{Ni (s)}
\end{equation}

En el ánodo, el metal se oxida y repone los iones metálicos en la solución, manteniendo constante la concentración de la sal:
\begin{equation}
    \text{Ni (s)} \rightarrow \text{Ni}^{2+} + 2e^-
\end{equation}

\subsection{Enjuague final y secado}

Se retira la pieza, se enjuaga y se seca, quedando con una capa metálica uniforme, brillante y resistente a la corrosión.


\end{document} 